\documentclass[11pt]{article}
\usepackage{amsmath, amssymb, booktabs}
\allowdisplaybreaks
\begin{document}

\section{Network and input set}
The network $f:\mathbb{R} \to \mathbb{R}$ with two hidden ReLU layers:
\begin{align*}
f(x) &= W^{(3)}\,\mathrm{ReLU}\bigl(W^{(2)}\,\mathrm{ReLU}\bigl(W^{(1)}\,x + b^{(1)}\bigr) + b^{(2)}\bigr) + b^{(3)}
\end{align*}
with
\begin{align*}
W^{(1)} &= \begin{bmatrix}
\frac{-3}{4} \\
-14
\end{bmatrix}, & b^{(1)} &= \bigl(\frac{1}{4},\; 2\bigr),
\end{align*}
\begin{align*}
W^{(2)} &= \begin{bmatrix}
-48 & \frac{7}{4} \\
-28 & \frac{7}{4}
\end{bmatrix}, & b^{(2)} &= \bigl(24,\; -4\bigr),
\end{align*}
\begin{align*}
W^{(3)} &= \begin{bmatrix}
\frac{-1}{2} & \frac{3}{2}
\end{bmatrix}, & b^{(3)} &= \bigl(7\bigr),
\end{align*}
and the input set
\begin{equation*}
 x \in [\underline{x}, \overline{x}] = [-1,\; 1].
\end{equation*}

\section{(i) Interval propagation}
Interval propagation (Box) with $W^{(k)+} = \max(W^{(k)}, 0)$ and
$W^{(k)-} = \min(W^{(k)}, 0)$, starting from $\underline{h}^{(0)} = \underline{x}$,
$\overline{h}^{(0)} = \overline{x}$:
\begin{align*}
\underline{z}^{(k)} &= W^{(k)+}\, \underline{h}^{(k-1)} + W^{(k)-}\, \overline{h}^{(k-1)} + b^{(k)}, &
\overline{z}^{(k)} &= W^{(k)+}\, \overline{h}^{(k-1)} + W^{(k)-}\, \underline{h}^{(k-1)} + b^{(k)},
\underline{h}^{(k)} &= \max(\underline{z}^{(k)}, 0), & \overline{h}^{(k)} &= \max(\overline{z}^{(k)}, 0) \quad (\text{ReLU layers}).
\end{align*}
Layer 1:
\begin{align*}
\underline{z^{(1)}} &= \bigl(\frac{-1}{2},\; -12\bigr)^{\!\top}, & \overline{z^{(1)}} &= \bigl(1,\; 16\bigr)^{\!\top}
 \\
\underline{h^{(1)}} &= \bigl(0,\; 0\bigr)^{\!\top}, & \overline{h^{(1)}} &= \bigl(1,\; 16\bigr)^{\!\top}
\end{align*}
Layer 2:
\begin{align*}
\underline{z^{(2)}} &= \bigl(-24,\; -32\bigr)^{\!\top}, & \overline{z^{(2)}} &= \bigl(52,\; 24\bigr)^{\!\top}
 \\
\underline{h^{(2)}} &= \bigl(0,\; 0\bigr)^{\!\top}, & \overline{h^{(2)}} &= \bigl(52,\; 24\bigr)^{\!\top}
\end{align*}
Layer 3:
\begin{align*}
\underline{z^{(3)}} &= \bigl(-19\bigr)^{\!\top}, & \overline{z^{(3)}} &= \bigl(43\bigr)^{\!\top}
\end{align*}
The interval upper bound on the output is $\overline{y}^{\text{interval}} = 43$.

\section{(ii) Symbolic interval propagation (shared relaxation domain)}
DPNFV with method \texttt{OuterBoundsNeurify} and max\_vars = 0. Each neuron carries symbolic lower and upper bound functions
$\underline{z}^{(k)}_i(\cdot) \le z^{(k)}_i \le \overline{z}^{(k)}_i(\cdot)$ (before activation)
and $\underline{h}^{(k)}_i(\cdot) \le h^{(k)}_i \le \overline{h}^{(k)}_i(\cdot)$ (after activation),
affine in the current symbols (initially only $x$).
Both relaxation slopes use the *shared* (outer) domain
$[\underline{z}^{\ell}_i, \overline{z}^{\ell}_i] = [\underline{z}^{u}_i, \overline{z}^{u}_i] = [\underline{z}^{(k)}_i, \overline{z}^{(k)}_i]$
of the pre-activation bounds, with $\beta_i = -\underline{z}^{\ell}_i$.
\paragraph{Layer 1 (ReLU).}
Symbolic pre-activation bounds after the linear map:
\begin{align*}
\underline{z}^{(1)}_{1} &= -\frac{3}{4}\,x + \frac{1}{4} \\
\underline{z}^{(1)}_{2} &= -14\,x + 2
\end{align*}
\begin{align*}
\overline{z}^{(1)}_{1} &= -\frac{3}{4}\,x + \frac{1}{4} \\
\overline{z}^{(1)}_{2} &= -14\,x + 2
\end{align*}
Branch domains and relaxation parameters:
\begin{align*}
[\underline{z}^{\ell}_i, \overline{z}^{\ell}_i] &= [\bigl(\frac{-1}{2},\; -12\bigr),\; \bigl(1,\; 16\bigr)], & [\underline{z}^{u}_i, \overline{z}^{u}_i] &= [\bigl(\frac{-1}{2},\; -12\bigr),\; \bigl(1,\; 16\bigr)], \\
\lambda^{\ell}_i &= \bigl(\frac{2}{3},\; \frac{4}{7}\bigr), & \lambda^{u}_i &= \bigl(\frac{2}{3},\; \frac{4}{7}\bigr), & \beta_i &= \bigl(\frac{1}{2},\; 12\bigr).
\end{align*}
with the relaxation rule
$\underline{h}^{(k)}_i = \lambda^{\ell}_i\, \underline{z}^{(k)}_i$ and
$\overline{h}^{(k)}_i = \lambda^{u}_i\, \overline{z}^{(k)}_i + \lambda^{u}_i \max(\beta_i, 0)$,
$\lambda^{\ell/u}_i = \mathrm{chord}(\underline{z}_i^{\ell/u}, \overline{z}_i^{\ell/u})$,
\mathrm{chord}(l, u) =
\begin{cases}
0, & u \le 0,\\
1, & l \ge 0,\\
\dfrac{u}{u - l}, & \text{otherwise (crossing ReLU)}.
\end{cases}
Relaxed post-activation bounds:
\begin{align*}
\underline{h}^{(1)}_{1} &= -\frac{1}{2}\,x + \frac{1}{6} \\
\underline{h}^{(1)}_{2} &= -8\,x + \frac{8}{7}
\end{align*}
\begin{align*}
\overline{h}^{(1)}_{1} &= -\frac{1}{2}\,x + \frac{1}{2} \\
\overline{h}^{(1)}_{2} &= -8\,x + 8
\end{align*}
\paragraph{Layer 2 (ReLU).}
Symbolic pre-activation bounds after the linear map:
\begin{align*}
\underline{z}^{(2)}_{1} &= 10\,x + 2 \\
\underline{z}^{(2)}_{2} &= -16
\end{align*}
\begin{align*}
\overline{z}^{(2)}_{1} &= 10\,x + 30 \\
\overline{z}^{(2)}_{2} &= \frac{16}{3}
\end{align*}
Branch domains and relaxation parameters:
\begin{align*}
[\underline{z}^{\ell}_i, \overline{z}^{\ell}_i] &= [\bigl(-8,\; -16\bigr),\; \bigl(12,\; -16\bigr)], & [\underline{z}^{u}_i, \overline{z}^{u}_i] &= [\bigl(20,\; \frac{16}{3}\bigr),\; \bigl(40,\; \frac{16}{3}\bigr)], \\
\lambda^{\ell}_i &= \bigl(\frac{5}{6},\; \frac{1}{4}\bigr), & \lambda^{u}_i &= \bigl(\frac{5}{6},\; \frac{1}{4}\bigr), & \beta_i &= \bigl(8,\; 16\bigr).
\end{align*}
with the relaxation rule
$\underline{h}^{(k)}_i = \lambda^{\ell}_i\, \underline{z}^{(k)}_i$ and
$\overline{h}^{(k)}_i = \lambda^{u}_i\, \overline{z}^{(k)}_i + \lambda^{u}_i \max(\beta_i, 0)$,
$\lambda^{\ell/u}_i = \mathrm{chord}(\underline{z}_i^{\ell/u}, \overline{z}_i^{\ell/u})$,
\mathrm{chord}(l, u) =
\begin{cases}
0, & u \le 0,\\
1, & l \ge 0,\\
\dfrac{u}{u - l}, & \text{otherwise (crossing ReLU)}.
\end{cases}
Relaxed post-activation bounds:
\begin{align*}
\underline{h}^{(2)}_{1} &= \frac{25}{3}\,x + \frac{5}{3} \\
\underline{h}^{(2)}_{2} &= -4
\end{align*}
\begin{align*}
\overline{h}^{(2)}_{1} &= \frac{25}{3}\,x + \frac{95}{3} \\
\overline{h}^{(2)}_{2} &= \frac{16}{3}
\end{align*}
\paragraph{Layer 3 (linear).}
Symbolic bounds of the output (substituted):
\begin{align*}
\underline{y}_{1} &= -\frac{25}{6}\,x - \frac{89}{6}
\end{align*}
\begin{align*}
\overline{y}_{1} &= -\frac{25}{6}\,x + \frac{85}{6}
\end{align*}
\paragraph{Output bounds.}
Substituting all remaining symbols gives the affine output bounds
\begin{align*}
\underline{y}_{1} &= -\frac{25}{6}\,x - \frac{89}{6}
\end{align*}
\begin{align*}
\overline{y}_{1} &= -\frac{25}{6}\,x + \frac{85}{6}
\end{align*}
\begin{equation*}
\underline{y}^{\text{DP}} = -19, \qquad \overline{y}^{\text{DP}} = \frac{55}{3}.
\end{equation*}

\section{(iii) Symbolic interval propagation (separate relaxation domains)}
DPNFV with method \texttt{NeurifyRelax} and max\_vars = 0. Each neuron carries symbolic lower and upper bound functions
$\underline{z}^{(k)}_i(\cdot) \le z^{(k)}_i \le \overline{z}^{(k)}_i(\cdot)$ (before activation)
and $\underline{h}^{(k)}_i(\cdot) \le h^{(k)}_i \le \overline{h}^{(k)}_i(\cdot)$ (after activation),
affine in the current symbols (initially only $x$).
The lower and the upper relaxation use *separate* domains:
the lower slope is computed on the range of the lower branch
$[\underline{z}^{\ell}_i, \overline{z}^{\ell}_i]$ (upper end from the upper-style substitution),
the upper slope on the range of the upper branch
$[\underline{z}^{u}_i, \overline{z}^{u}_i]$ (lower end from the lower-style substitution),
with $\beta_i = -\underline{z}^{u}_i$.
\paragraph{Layer 1 (ReLU).}
Symbolic pre-activation bounds after the linear map:
\begin{align*}
\underline{z}^{(1)}_{1} &= -\frac{3}{4}\,x + \frac{1}{4} \\
\underline{z}^{(1)}_{2} &= -14\,x + 2
\end{align*}
\begin{align*}
\overline{z}^{(1)}_{1} &= -\frac{3}{4}\,x + \frac{1}{4} \\
\overline{z}^{(1)}_{2} &= -14\,x + 2
\end{align*}
Branch domains and relaxation parameters:
\begin{align*}
[\underline{z}^{\ell}_i, \overline{z}^{\ell}_i] &= [\bigl(\frac{-1}{2},\; -12\bigr),\; \bigl(1,\; 16\bigr)], & [\underline{z}^{u}_i, \overline{z}^{u}_i] &= [\bigl(\frac{-1}{2},\; -12\bigr),\; \bigl(1,\; 16\bigr)], \\
\lambda^{\ell}_i &= \bigl(\frac{2}{3},\; \frac{4}{7}\bigr), & \lambda^{u}_i &= \bigl(\frac{2}{3},\; \frac{4}{7}\bigr), & \beta_i &= \bigl(\frac{1}{2},\; 12\bigr).
\end{align*}
with the relaxation rule
$\underline{h}^{(k)}_i = \lambda^{\ell}_i\, \underline{z}^{(k)}_i$ and
$\overline{h}^{(k)}_i = \lambda^{u}_i\, \overline{z}^{(k)}_i + \lambda^{u}_i \max(\beta_i, 0)$,
$\lambda^{\ell/u}_i = \mathrm{chord}(\underline{z}_i^{\ell/u}, \overline{z}_i^{\ell/u})$,
\mathrm{chord}(l, u) =
\begin{cases}
0, & u \le 0,\\
1, & l \ge 0,\\
\dfrac{u}{u - l}, & \text{otherwise (crossing ReLU)}.
\end{cases}
Relaxed post-activation bounds:
\begin{align*}
\underline{h}^{(1)}_{1} &= -\frac{1}{2}\,x + \frac{1}{6} \\
\underline{h}^{(1)}_{2} &= -8\,x + \frac{8}{7}
\end{align*}
\begin{align*}
\overline{h}^{(1)}_{1} &= -\frac{1}{2}\,x + \frac{1}{2} \\
\overline{h}^{(1)}_{2} &= -8\,x + 8
\end{align*}
\paragraph{Layer 2 (ReLU).}
Symbolic pre-activation bounds after the linear map:
\begin{align*}
\underline{z}^{(2)}_{1} &= 10\,x + 2 \\
\underline{z}^{(2)}_{2} &= -16
\end{align*}
\begin{align*}
\overline{z}^{(2)}_{1} &= 10\,x + 30 \\
\overline{z}^{(2)}_{2} &= \frac{16}{3}
\end{align*}
Branch domains and relaxation parameters:
\begin{align*}
[\underline{z}^{\ell}_i, \overline{z}^{\ell}_i] &= [\bigl(-8,\; -16\bigr),\; \bigl(12,\; -16\bigr)], & [\underline{z}^{u}_i, \overline{z}^{u}_i] &= [\bigl(20,\; \frac{16}{3}\bigr),\; \bigl(40,\; \frac{16}{3}\bigr)], \\
\lambda^{\ell}_i &= \bigl(\frac{3}{5},\; 0\bigr), & \lambda^{u}_i &= \bigl(1,\; 1\bigr), & \beta_i &= \bigl(-20,\; \frac{-16}{3}\bigr).
\end{align*}
with the relaxation rule
$\underline{h}^{(k)}_i = \lambda^{\ell}_i\, \underline{z}^{(k)}_i$ and
$\overline{h}^{(k)}_i = \lambda^{u}_i\, \overline{z}^{(k)}_i + \lambda^{u}_i \max(\beta_i, 0)$,
$\lambda^{\ell/u}_i = \mathrm{chord}(\underline{z}_i^{\ell/u}, \overline{z}_i^{\ell/u})$,
\mathrm{chord}(l, u) =
\begin{cases}
0, & u \le 0,\\
1, & l \ge 0,\\
\dfrac{u}{u - l}, & \text{otherwise (crossing ReLU)}.
\end{cases}
Relaxed post-activation bounds:
\begin{align*}
\underline{h}^{(2)}_{1} &= 6\,x + \frac{6}{5} \\
\underline{h}^{(2)}_{2} &= 0
\end{align*}
\begin{align*}
\overline{h}^{(2)}_{1} &= 10\,x + 30 \\
\overline{h}^{(2)}_{2} &= \frac{16}{3}
\end{align*}
\paragraph{Layer 3 (linear).}
Symbolic bounds of the output (substituted):
\begin{align*}
\underline{y}_{1} &= -5\,x - 8
\end{align*}
\begin{align*}
\overline{y}_{1} &= -3\,x + \frac{72}{5}
\end{align*}
\paragraph{Output bounds.}
Substituting all remaining symbols gives the affine output bounds
\begin{align*}
\underline{y}_{1} &= -5\,x - 8
\end{align*}
\begin{align*}
\overline{y}_{1} &= -3\,x + \frac{72}{5}
\end{align*}
\begin{equation*}
\underline{y}^{\text{DP}} = -13, \qquad \overline{y}^{\text{DP}} = \frac{87}{5}.
\end{equation*}

\section{(iii') Symbolic interval propagation with fresh variables}
DPNFV with method \texttt{NeurifyRelax} and max\_vars = 10. Each neuron carries symbolic lower and upper bound functions
$\underline{z}^{(k)}_i(\cdot) \le z^{(k)}_i \le \overline{z}^{(k)}_i(\cdot)$ (before activation)
and $\underline{h}^{(k)}_i(\cdot) \le h^{(k)}_i \le \overline{h}^{(k)}_i(\cdot)$ (after activation),
affine in the current symbols (initially only $x$).
As in (iii), but the first neuron chosen below keeps its *exact*
symbolic dependence as a fresh variable $\xi$: instead of relaxing it right away, $\xi = h^{(k)}_v$
is carried symbolically and only substituted at the next layer. The stored symbolic bounds of $\xi$
are nevertheless already \emph{relaxed} ones: they are the chord relaxations of the neuron, i.e.
the substituted bounds of $z^{(k)}_v$ multiplied by the neuron's slopes (and shifted by the upper
chord offset), so they inherit exactly the relaxation that rung (iii) would have applied.
\paragraph{Layer 1 (ReLU).}
Symbolic pre-activation bounds after the linear map:
\begin{align*}
\underline{z}^{(1)}_{1} &= -\frac{3}{4}\,x + \frac{1}{4} \\
\underline{z}^{(1)}_{2} &= -14\,x + 2
\end{align*}
\begin{align*}
\overline{z}^{(1)}_{1} &= -\frac{3}{4}\,x + \frac{1}{4} \\
\overline{z}^{(1)}_{2} &= -14\,x + 2
\end{align*}
The crossing ReLU with the largest range is neuron 2; it becomes a fresh variable. Its symbolic bounds below are the already relaxed (chord) bounds of the neuron:
\begin{align*}
\xi &= h^{(1)}_{2}, & \underline{\xi}(x) &= -8\,x + \frac{8}{7}, & \overline{\xi}(x) &= -8\,x + 8.
\end{align*}
Branch domains and relaxation parameters:
\begin{align*}
[\underline{z}^{\ell}_i, \overline{z}^{\ell}_i] &= [\bigl(\frac{-1}{2},\; -12\bigr),\; \bigl(1,\; 16\bigr)], & [\underline{z}^{u}_i, \overline{z}^{u}_i] &= [\bigl(\frac{-1}{2},\; -12\bigr),\; \bigl(1,\; 16\bigr)], \\
\lambda^{\ell}_i &= \bigl(\frac{2}{3},\; \frac{4}{7}\bigr), & \lambda^{u}_i &= \bigl(\frac{2}{3},\; \frac{4}{7}\bigr), & \beta_i &= \bigl(\frac{1}{2},\; 12\bigr).
\end{align*}
with the relaxation rule
$\underline{h}^{(k)}_i = \lambda^{\ell}_i\, \underline{z}^{(k)}_i$ and
$\overline{h}^{(k)}_i = \lambda^{u}_i\, \overline{z}^{(k)}_i + \lambda^{u}_i \max(\beta_i, 0)$,
$\lambda^{\ell/u}_i = \mathrm{chord}(\underline{z}_i^{\ell/u}, \overline{z}_i^{\ell/u})$,
\mathrm{chord}(l, u) =
\begin{cases}
0, & u \le 0,\\
1, & l \ge 0,\\
\dfrac{u}{u - l}, & \text{otherwise (crossing ReLU)}.
\end{cases}
Relaxed post-activation bounds:
\begin{align*}
\underline{h}^{(1)}_{1} &= -\frac{1}{2}\,x + \frac{1}{6} \\
\underline{h}^{(1)}_{2} &= x
\end{align*}
\begin{align*}
\overline{h}^{(1)}_{1} &= -\frac{1}{2}\,x + \frac{1}{2} \\
\overline{h}^{(1)}_{2} &= x
\end{align*}
\paragraph{Layer 2 (ReLU).}
Symbolic pre-activation bounds after the linear map:
\begin{align*}
\underline{z}^{(2)}_{1} &= 24\,x + \frac{7}{4}\,\xi \\
\underline{z}^{(2)}_{2} &= 14\,x + \frac{7}{4}\,\xi - 18
\end{align*}
\begin{align*}
\overline{z}^{(2)}_{1} &= 24\,x + \frac{7}{4}\,\xi + 16 \\
\overline{z}^{(2)}_{2} &= 14\,x + \frac{7}{4}\,\xi - \frac{26}{3}
\end{align*}
Substituting $\xi$ by its symbolic bounds ($c\,\xi \mapsto c^{+}\,\underline{\xi}(x) + c^{-}\,\overline{\xi}(x)$ in lower-style, and crossed in upper-style substitutions):
\begin{align*}
\underline{z}^{(2),\ell}_{1} &= 10\,x + 2 \\
\underline{z}^{(2),\ell}_{2} &= -16
\end{align*}
\begin{align*}
\overline{z}^{(2),\ell}_{1} &= 10\,x + 14 \\
\overline{z}^{(2),\ell}_{2} &= -4
\end{align*}
\begin{align*}
\underline{z}^{(2),u}_{1} &= 10\,x + 18 \\
\underline{z}^{(2),u}_{2} &= -\frac{20}{3}
\end{align*}
\begin{align*}
\overline{z}^{(2),u}_{1} &= 10\,x + 30 \\
\overline{z}^{(2),u}_{2} &= \frac{16}{3}
\end{align*}
Branch domains and relaxation parameters:
\begin{align*}
[\underline{z}^{\ell}_i, \overline{z}^{\ell}_i] &= [\bigl(-8,\; -16\bigr),\; \bigl(24,\; -4\bigr)], & [\underline{z}^{u}_i, \overline{z}^{u}_i] &= [\bigl(8,\; \frac{-20}{3}\bigr),\; \bigl(40,\; \frac{16}{3}\bigr)], \\
\lambda^{\ell}_i &= \bigl(\frac{3}{4},\; 0\bigr), & \lambda^{u}_i &= \bigl(1,\; \frac{4}{9}\bigr), & \beta_i &= \bigl(-8,\; \frac{20}{3}\bigr).
\end{align*}
with the relaxation rule
$\underline{h}^{(k)}_i = \lambda^{\ell}_i\, \underline{z}^{(k)}_i$ and
$\overline{h}^{(k)}_i = \lambda^{u}_i\, \overline{z}^{(k)}_i + \lambda^{u}_i \max(\beta_i, 0)$,
$\lambda^{\ell/u}_i = \mathrm{chord}(\underline{z}_i^{\ell/u}, \overline{z}_i^{\ell/u})$,
\mathrm{chord}(l, u) =
\begin{cases}
0, & u \le 0,\\
1, & l \ge 0,\\
\dfrac{u}{u - l}, & \text{otherwise (crossing ReLU)}.
\end{cases}
Relaxed post-activation bounds:
\begin{align*}
\underline{h}^{(2)}_{1} &= 18\,x + \frac{21}{16}\,\xi \\
\underline{h}^{(2)}_{2} &= 0
\end{align*}
\begin{align*}
\overline{h}^{(2)}_{1} &= 24\,x + \frac{7}{4}\,\xi + 16 \\
\overline{h}^{(2)}_{2} &= \frac{56}{9}\,x + \frac{7}{9}\,\xi - \frac{8}{9}
\end{align*}
\paragraph{Layer 3 (linear).}
Symbolic bounds of the output (substituted):
\begin{align*}
\underline{y}_{1} &= -5\,x - 8
\end{align*}
\begin{align*}
\overline{y}_{1} &= -\frac{15}{4}\,x + \frac{39}{4}
\end{align*}
\paragraph{Output bounds.}
Substituting all remaining symbols gives the affine output bounds
\begin{align*}
\underline{y}_{1} &= -5\,x - 8
\end{align*}
\begin{align*}
\overline{y}_{1} &= -\frac{15}{4}\,x + \frac{39}{4}
\end{align*}
\begin{equation*}
\underline{y}^{\text{DP}} = -13, \qquad \overline{y}^{\text{DP}} = \frac{27}{2}.
\end{equation*}

\section{(iv) Zonotope propagation (Theorem 3.1, DeepZ)}
A zonotope is $Z = \langle c, G \rangle = \{c + G\epsilon : \epsilon \in [-1,1]^m\}$
(center $c$, generators the columns of $G$); its exact interval bounds are
$c \pm \sum_j |G_{:j}|$. The input set becomes $Z^{(0)} = \langle c^{(0)}, G^{(0)} \rangle$ with
$c^{(0)} = (\underline{x} + \overline{x})/2$ and $G^{(0)}$ the diagonal of $(\overline{x} - \underline{x})/2$.
Per ReLU neuron with interval bounds $[l_i, u_i]$:
$\lambda_i = \mathrm{chord}(l_i, u_i)$ and, for crossing neurons,
$\mathrm{ReLU}(z_i) \le \lambda_i (z_i - l_i) = \lambda_i z_i + \mu_i$ with error-generator magnitude
$\mu_i = -\lambda_i l_i / 2 > 0$. The rectified zonotope is
\begin{equation*}
c' = \lambda \odot c + \mu, \qquad G' = \bigl[\, \lambda \odot G \;\big|\; \mu\, \bigr]
\end{equation*}
($\lambda_i = 1, \mu_i = 0$ for active, $\lambda_i = 0$ for inactive neurons;
the new error generator $\mu_i$ has entries only in row $i$).
For the example network:
\paragraph{Layer 1.}
Pre-activation zonotope: $c = \bigl(\frac{1}{4},\; 2\bigr)$, generator columns:
\begin{align*}
g_{1} &= \bigl(\frac{-3}{4},\; -14\bigr) \quad (\epsilon_{1}).
\end{align*}
Interval bounds: $l = \bigl(\frac{-1}{2},\; -12\bigr), \; u = \bigl(1,\; 16\bigr)$. 
Crossing ReLUs at neurons 1, 2 with
\begin{align*}
\lambda &= \bigl(\frac{2}{3},\; \frac{4}{7}\bigr), & \mu &= \bigl(\frac{1}{6},\; \frac{24}{7}\bigr).
\end{align*}
Post-activation zonotope: $c' = \bigl(\frac{1}{3},\; \frac{32}{7}\bigr)$, generators:
\begin{align*}
g'_{1} &= \bigl(\frac{-1}{2},\; -8\bigr), &
g'_{2} &= \bigl(\frac{1}{6},\; 0\bigr), &
g'_{3} &= \bigl(0,\; \frac{24}{7}\bigr).
\end{align*}
\paragraph{Layer 2.}
Pre-activation zonotope: $c = \bigl(16,\; \frac{-16}{3}\bigr)$, generator columns:
\begin{align*}
g_{1} &= \bigl(10,\; 0\bigr) \quad (\epsilon_{1}), &
g_{2} &= \bigl(-8,\; \frac{-14}{3}\bigr) \quad (\eta^{(1)}_{1}), &
g_{3} &= \bigl(6,\; 6\bigr) \quad (\eta^{(1)}_{2}).
\end{align*}
Interval bounds: $l = \bigl(-8,\; -16\bigr), \; u = \bigl(40,\; \frac{16}{3}\bigr)$. 
Crossing ReLUs at neurons 1, 2 with
\begin{align*}
\lambda &= \bigl(\frac{5}{6},\; \frac{1}{4}\bigr), & \mu &= \bigl(\frac{10}{3},\; 2\bigr).
\end{align*}
Post-activation zonotope: $c' = \bigl(\frac{50}{3},\; \frac{2}{3}\bigr)$, generators:
\begin{align*}
g'_{1} &= \bigl(\frac{25}{3},\; 0\bigr), &
g'_{2} &= \bigl(\frac{-20}{3},\; \frac{-7}{6}\bigr), &
g'_{3} &= \bigl(5,\; \frac{3}{2}\bigr), &
g'_{4} &= \bigl(\frac{10}{3},\; 0\bigr), &
g'_{5} &= \bigl(0,\; 2\bigr).
\end{align*}
\paragraph{Layer 3.}
Pre-activation zonotope: $c = \bigl(\frac{-1}{3}\bigr)$, generator columns:
\begin{align*}
g_{1} &= \bigl(\frac{32}{3}\bigr) \quad (\epsilon_{1}).
\end{align*}
Interval bounds: $l = \bigl(-11\bigr), \; u = \bigl(\frac{31}{3}\bigr)$. 
No crossing ReLUs ($\lambda = \bigl(1\bigr)$).
Post-activation zonotope: $c' = \bigl(\frac{-1}{3}\bigr)$, generators:
\begin{align*}
g'_{1} &= \bigl(\frac{32}{3}\bigr).
\end{align*}
The zonotope upper bound on the output is $\overline{y}^{\text{zono}} = c + \sum_j |g_j| = \frac{31}{3}$.

\section{(v) CROWN backsubstitution (aCROWN initialisation pass)}
CROWN computes the output bound by \emph{backsubstitution}: affine bound functions
$\underline{g}(h) = \Lambda h + \gamma \le z$ and $\overline{g}(h) = \Gamma h + \gamma \ge z$
are propagated backwards layer by layer, relaxing each ReLU by chord bounds whose slopes are taken
from the concrete pre-activation bounds found in earlier stages of the same pass. For a neuron with
relaxation domain $[\ell_i, u_i]$ (the concrete bounds of $z^{(k)}_i$) define
\begin{align*}
a^{\ell}_i &= \begin{cases} 1, & \ell_i \ge 0,\\ 0, & \text{otherwise}, \end{cases} & a^{u}_i &= \mathrm{chord}(\ell_i, u_i), & b^{u}_i &= a^{u}_i \max(-\ell_i, 0),
\end{align*}
so that $\mathrm{ReLU}(z_i) \ge a^{\ell}_i z_i$ and
$\mathrm{ReLU}(z_i) \le a^{u}_i z_i + b^{u}_i$. A bound in terms of the post-activations $h$ becomes
one in terms of the pre-activations $z$ via
\begin{align*}
\Lambda' &= \Lambda^{+} a^{\ell} + \Lambda^{-} a^{u}, & \gamma' &= \Lambda^{-} b^{u} + \gamma && \text{(lower)}, \
\Gamma' &= \Gamma^{+} a^{u} + \Gamma^{-} a^{\ell}, & \gamma' &= \Gamma^{+} b^{u} + \gamma && \text{(upper)},
\end{align*}
with $\Lambda^{+} = \max(\Lambda, 0)$, $\Lambda^{-} = \min(\Lambda, 0)$ (analogously $\Gamma$) and the
slope vectors acting column-wise (diagonal scaling). Composing with
$z^{(k)} = W^{(k)} h^{(k-1)} + b^{(k)}$ yields the bound in terms of $h^{(k-1)}$. The pass runs in
stages $i = 1, \dots, n$: stage $i$ backsubstitutes the pre-activation $z^{(i)}$ to the input, using
the concrete bounds recorded by stages $< i$ as relaxation domains, and records
$[\underline{z}^{(i)}, \overline{z}^{(i)}]$ for the later stages. Within each stage, $g$ and $g'$
denote the current lower/upper bound functions (redefined at every step).
\paragraph{Stage 1: bounds on $z^{(1)}$.}
The bound functions are the layer map $z^{(1)} = W^{(1)} x + b^{(1)}$ itself; both branches coincide and their interval evaluation is the interval bound of rung (i).
Interval evaluation of both branches over $[\underline{x}, \overline{x}]$ gives
\begin{align*}
[\underline{z}^{\ell}, \overline{z}^{\ell}] &= [\bigl(\frac{-1}{2},\; -12\bigr),\; \bigl(1,\; 16\bigr)] \quad\text{(lower branch)}, & [\underline{z}^{u}, \overline{z}^{u}] &= [\bigl(\frac{-1}{2},\; -12\bigr),\; \bigl(1,\; 16\bigr)] \quad\text{(upper branch)}, \\
\text{recorded: } [\underline{z^{(1)}}, \overline{z^{(1)}}] &= [\bigl(\frac{-1}{2},\; -12\bigr),\; \bigl(1,\; 16\bigr)] \quad\text{(lower end from the lower, upper end from the upper branch)}.
\end{align*}
\paragraph{Stage 2: bounds on $z^{(2)}$.}
Backsubstituting through layer 1 (ReLU) with relaxation domain $[\underline{z}^{(j)}, \overline{z}^{(j)}] = [\bigl(\frac{-1}{2},\; -12\bigr),\; \bigl(1,\; 16\bigr)]$ (the stage-1 result):
\begin{align*}
a^{\ell} &= \bigl(1,\; 1\bigr), & a^{u} &= \bigl(\frac{2}{3},\; \frac{4}{7}\bigr), & b^{u} &= \bigl(\frac{1}{3},\; \frac{48}{7}\bigr).
\end{align*}
Applying the relaxation rules gives bound functions in terms of $z^{(1)}$:
\begin{align*}
\underline{g}_{1} &= -32\,z^{(1)}_{1} + \frac{7}{4}\,z^{(1)}_{2} + 8 \\
\underline{g}_{2} &= -\frac{56}{3}\,z^{(1)}_{1} + \frac{7}{4}\,z^{(1)}_{2} - \frac{40}{3}
\end{align*}
\begin{align*}
\overline{g}_{1} &= -48\,z^{(1)}_{1} + z^{(1)}_{2} + 36 \\
\overline{g}_{2} &= -28\,z^{(1)}_{1} + z^{(1)}_{2} + 8
\end{align*}
Composing with $z^{(1)} = W^{(1)} x + b^{(1)}$ gives bound functions in terms of $x$:
\begin{align*}
\underline{g}'_{1} &= -\frac{1}{2}\,x + \frac{7}{2} \\
\underline{g}'_{2} &= -\frac{21}{2}\,x - \frac{29}{2}
\end{align*}
\begin{align*}
\overline{g}'_{1} &= 22\,x + 26 \\
\overline{g}'_{2} &= 7\,x + 3
\end{align*}
Interval evaluation of both branches over $[\underline{x}, \overline{x}]$ gives
\begin{align*}
[\underline{z}^{\ell}, \overline{z}^{\ell}] &= [\bigl(3,\; -25\bigr),\; \bigl(4,\; -4\bigr)] \quad\text{(lower branch)}, & [\underline{z}^{u}, \overline{z}^{u}] &= [\bigl(4,\; -4\bigr),\; \bigl(48,\; 10\bigr)] \quad\text{(upper branch)}, \\
\text{recorded: } [\underline{z^{(2)}}, \overline{z^{(2)}}] &= [\bigl(3,\; -25\bigr),\; \bigl(48,\; 10\bigr)] \quad\text{(lower end from the lower, upper end from the upper branch)}.
\end{align*}
\paragraph{Stage 3: bounds on the output $y$.}
Backsubstituting through layer 2 (ReLU) with relaxation domain $[\underline{z}^{(j)}, \overline{z}^{(j)}] = [\bigl(3,\; -25\bigr),\; \bigl(48,\; 10\bigr)]$ (the stage-2 result):
\begin{align*}
a^{\ell} &= \bigl(1,\; 0\bigr), & a^{u} &= \bigl(1,\; \frac{2}{7}\bigr), & b^{u} &= \bigl(0,\; \frac{50}{7}\bigr).
\end{align*}
Applying the relaxation rules gives bound functions in terms of $z^{(2)}$:
\begin{align*}
\underline{g}_{1} &= -\frac{1}{2}\,z^{(2)}_{1} + 7
\end{align*}
\begin{align*}
\overline{g}_{1} &= -\frac{1}{2}\,z^{(2)}_{1} + \frac{3}{7}\,z^{(2)}_{2} + \frac{124}{7}
\end{align*}
Composing with $z^{(j)} = W^{(j)} h^{1} + b^{(j)}$ gives bound functions in terms of $h^{1}$:
\begin{align*}
\underline{g}'_{1} &= 24\,h^{1}_{1} - \frac{7}{8}\,h^{1}_{2} - 5
\end{align*}
\begin{align*}
\overline{g}'_{1} &= 12\,h^{1}_{1} - \frac{1}{8}\,h^{1}_{2} + 4
\end{align*}
Backsubstituting through layer 1 (ReLU) with relaxation domain $[\underline{z}^{(j)}, \overline{z}^{(j)}] = [\bigl(\frac{-1}{2},\; -12\bigr),\; \bigl(1,\; 16\bigr)]$ (the stage-1 result):
\begin{align*}
a^{\ell} &= \bigl(1,\; 1\bigr), & a^{u} &= \bigl(\frac{2}{3},\; \frac{4}{7}\bigr), & b^{u} &= \bigl(\frac{1}{3},\; \frac{48}{7}\bigr).
\end{align*}
Applying the relaxation rules gives bound functions in terms of $z^{(1)}$:
\begin{align*}
\underline{g}_{1} &= 24\,z^{(1)}_{1} - \frac{1}{2}\,z^{(1)}_{2} - 11
\end{align*}
\begin{align*}
\overline{g}_{1} &= 8\,z^{(1)}_{1} - \frac{1}{8}\,z^{(1)}_{2} + 8
\end{align*}
Composing with $z^{(1)} = W^{(1)} x + b^{(1)}$ gives bound functions in terms of $x$:
\begin{align*}
\underline{g}'_{1} &= -11\,x - 6
\end{align*}
\begin{align*}
\overline{g}'_{1} &= -\frac{17}{4}\,x + \frac{39}{4}
\end{align*}
Interval evaluation of both branches over $[\underline{x}, \overline{x}]$ gives
\begin{align*}
[\underline{z}^{\ell}, \overline{z}^{\ell}] &= [\bigl(-17\bigr),\; \bigl(5\bigr)] \quad\text{(lower branch)}, & [\underline{z}^{u}, \overline{z}^{u}] &= [\bigl(\frac{11}{2}\bigr),\; \bigl(14\bigr)] \quad\text{(upper branch)}, \\
\text{recorded: } [\underline{y}, \overline{y}] &= [\bigl(-17\bigr),\; \bigl(14\bigr)] \quad\text{(lower end from the lower, upper end from the upper branch)}.
\end{align*}
\paragraph{Output bound.}
\begin{equation*}
\underline{y}^{\text{CROWN}} = -17, \qquad \overline{y}^{\text{CROWN}} = 14.
\end{equation*}

\section{(vi) NNPoly with quadratic relaxations (PolyCROWN initialisation pass)}
With the single input $x = c + r\,\epsilon$, $\epsilon \in [-1, 1]$ (center $c = \bigl(0\bigr)$, radius $r = \bigl(1\bigr)$), the first-layer pre-activations are exact affine polynomials in $\epsilon$, $z^{(1)}_i = z^{(1)}_{i,0} + z^{(1)}_{i,1}\,\epsilon$:
\begin{align*}
z^{(1)}_{1} &= -\frac{3}{4}\,\epsilon + \frac{1}{4} \\
z^{(1)}_{2} &= -14\,\epsilon + 2
\end{align*}
Their exact bounds $\underline{z}^{(1)}_i = z^{(1)}_{i,0} - |z^{(1)}_{i,1}|$ and $\overline{z}^{(1)}_i = z^{(1)}_{i,0} + |z^{(1)}_{i,1}|$ agree with the interval rung (i):
\begin{align*}
\underline{z^{(1)}} &= \bigl(\frac{-1}{2},\; -12\bigr)^{\!\top}, & \overline{z^{(1)}} &= \bigl(1,\; 16\bigr)^{\!\top}.
\end{align*}
Each crossing ReLU neuron (with $\ell_i = \underline{z}^{(1)}_i < 0 < u_i = \overline{z}^{(1)}_i$)
receives a quadratic (CROWNQuad) relaxation; active and inactive neurons use $\mathrm{ReLU}(z) \le z$
and $\mathrm{ReLU}(z) = 0$. For crossing neurons:
\begin{align*}
\text{upper: } \mathrm{ReLU}(z) &\le \begin{cases} \dfrac{-l u^2 + (l^2 + u^2)\, z - l\, z^2}{(l - u)^2}, & -l \le u, \\[3mm] \dfrac{l^2 u - 2 l u\, z + u\, z^2}{(l - u)^2}, & -l > u, \end{cases} \\[2mm]
\text{lower: } \mathrm{ReLU}(z) &\ge \begin{cases} z, & u \ge -2l, \\[3mm] \dfrac{-l u\, z + u\, z^2}{u^2 - l u}, & -l < u < -2l, \\[3mm] 0, & u \le -l. \end{cases}
\end{align*}
Per-neuron monomial coefficients $c_i = (c_0, c_1, c_2)$ with $\mathrm{ReLU}(z_i) \lesseqgtr c_0 + c_1 z_i + c_2 z_i^2$:
\begin{align*}
c^{\ell}_{1} &= \bigl(0,\; 1,\; 0\bigr), & c^{u}_{1} &= \bigl(\frac{2}{9},\; \frac{5}{9},\; \frac{2}{9}\bigr) \\
c^{\ell}_{2} &= \bigl(0,\; \frac{3}{7},\; \frac{1}{28}\bigr), & c^{u}_{2} &= \bigl(\frac{192}{49},\; \frac{25}{49},\; \frac{3}{196}\bigr)
\end{align*}
Substituting $z^{(1)}_i(\epsilon) = z^{(1)}_{i,0} + z^{(1)}_{i,1} \epsilon$ into the relaxation
polynomials (with $z^2 = z_0^2 + 2 z_0 z_1 \epsilon + z_1^2 \epsilon^2$) gives the quolynomial
(degree-2 polynomial) bounds of the post-activations:
\begin{align*}
\hat{h}^{(1),\ell}_{1} &= -\frac{3}{4}\,\epsilon + \frac{1}{4} \\
\hat{h}^{(1),\ell}_{2} &= -8\,\epsilon + 7\,\epsilon^2 + 1
\end{align*}
\begin{align*}
\hat{h}^{(1),u}_{1} &= -\frac{1}{2}\,\epsilon + \frac{1}{8}\,\epsilon^2 + \frac{3}{8} \\
\hat{h}^{(1),u}_{2} &= -8\,\epsilon + 3\,\epsilon^2 + 5
\end{align*}
Concrete bounds of a quolynomial $q(\epsilon) = q_0 + q_1 \epsilon + q_2 \epsilon^2$ over
$\epsilon \in [-1, 1]$ follow from the monomial ranges $1 \in [1, 1]$, $\epsilon \in [-1, 1]$,
$\epsilon^2 \in [0, 1]$, i.e. from $q_0 \pm |q_1|$ and the sign of $q_2$:
\begin{equation*}
q \in \bigl[\, q_0 - |q_1| + \min(q_2, 0), \;\; q_0 + |q_1| + \max(q_2, 0) \,\bigr].
\end{equation*}
Applied to the quolynomial bounds of $h^{(1)}$ (lower branch for the lower end, upper branch for the upper end):
\begin{align*}
\underline{\hat{h}^{(1)}} &= \bigl(\frac{-1}{2},\; -7\bigr)^{\!\top}, & \overline{\hat{h}^{(1)}} &= \bigl(1,\; 16\bigr)^{\!\top}.
\end{align*}
Composing the linear map $z^{(2)} = W^{(2)} h^{(1)} + b^{(2)}$ with the quolynomial bounds
(positive coefficients read the lower bound polynomial $\hat{h}^{\ell}$, negative coefficients the
upper one $\hat{h}^{u}$; the bias enters the constant monomial) gives quolynomial bounds of
$z^{(2)}$:
\begin{align*}
\underline{z}^{(2)}_{1} &= 10\,\epsilon + \frac{25}{4}\,\epsilon^2 + \frac{31}{4} \\
\underline{z}^{(2)}_{2} &= \frac{35}{4}\,\epsilon^2 - \frac{51}{4}
\end{align*}
\begin{align*}
\overline{z}^{(2)}_{1} &= 22\,\epsilon + \frac{21}{4}\,\epsilon^2 + \frac{83}{4} \\
\overline{z}^{(2)}_{2} &= 7\,\epsilon + \frac{21}{4}\,\epsilon^2 - \frac{9}{4}
\end{align*}
Their concrete bounds (same monomial rule, now over the mph{quolynomial} domain, which is tighter than the interval rung):
\begin{align*}
\underline{z^{(2)}} &= \bigl(\frac{-9}{4},\; \frac{-51}{4}\bigr)^{\!\top}, & \overline{z^{(2)}} &= \bigl(48,\; 10\bigr)^{\!\top}.
\end{align*}
The output $y = W^{(3)} h^{(2)} + b^{(3)}$ is backsubstituted through ReLU layer 2 with the linear chord rules of rung (v), using the quolynomial-domain bounds
$[\underline{z}^{(2)}, \overline{z}^{(2)}]$ as relaxation domains:
\begin{align*}
a^{\ell} &= \bigl(1,\; 0\bigr), & a^{u} &= \bigl(\frac{64}{67},\; \frac{40}{91}\bigr), & b^{u} &= \bigl(\frac{144}{67},\; \frac{510}{91}\bigr).
\end{align*}
Relaxing ReLU layer 2 gives bound functions in terms of $z^{(2)}$:
\begin{align*}
\underline{g}_{1} &= -\frac{32}{67}\,z^{(2)}_{1} + \frac{397}{67}
\end{align*}
\begin{align*}
\overline{g}_{1} &= -\frac{1}{2}\,z^{(2)}_{1} + \frac{60}{91}\,z^{(2)}_{2} + \frac{1402}{91}
\end{align*}
Composing with $z^{(2)} = W^{(2)} h^{(1)} + b^{(2)}$ gives bound functions in terms of $h^{(1)}$:
\begin{align*}
\underline{g}'_{1} &= \frac{1536}{67}\,h^{(1)}_{1} - \frac{56}{67}\,h^{(1)}_{2} - \frac{371}{67}
\end{align*}
\begin{align*}
\overline{g}'_{1} &= \frac{72}{13}\,h^{(1)}_{1} + \frac{29}{104}\,h^{(1)}_{2} + \frac{10}{13}
\end{align*}
Composing these bound functions with the quolynomial bounds of $h^{(1)}$ (positive
coefficients read $\hat{h}^{u}$ resp.\ $\hat{h}^{\ell}$ as appropriate, the shift $\gamma$ enters the
constant monomial) yields polynomial bounds on the output:
\begin{align*}
\underline{y}_{1} &= -\frac{704}{67}\,\epsilon - \frac{168}{67}\,\epsilon^2 - \frac{267}{67}
\end{align*}
\begin{align*}
\overline{y}_{1} &= -5\,\epsilon + \frac{159}{104}\,\epsilon^2 + \frac{441}{104}
\end{align*}
Evaluating both over the monomial domain gives
\begin{equation*}
\underline{y}^{\text{NNPoly}} = -17, \qquad \overline{y}^{\text{NNPoly}} = \frac{140}{13}.
\end{equation*}

\section{Summary: the bound ladder}
Collecting the upper bounds; they are ordered loose to tight from (i) to (vi)
(and further tightened by the optimised variants):
\begin{table}[ht]\centering
\begin{tabular}{llc}\toprule
rung & algorithm & $\overline{y}$ \\
\midrule
(i) & interval propagation & 43 \\
(ii) & symbolic interval (shared domains) & \frac{55}{3} \\
(iii) & symbolic interval (separate domains) & \frac{87}{5} \\
(iii') & symbolic interval + fresh variables & \frac{27}{2} \\
(iv) & zonotope propagation (DeepZ) & \frac{31}{3} \\
(v) & CROWN backsubstitution & 14 \\
(vi) & NNPoly (quadratic relaxations) & \frac{140}{13} \\
\bottomrule
\end{tabular}
\end{table}


\end{document}
